\documentclass[11pt,a4paper]{article}

% ── Codificación y fuentes ────────────────────────────────────────────────────
\usepackage[utf8]{inputenc}
\usepackage[T1]{fontenc}
\usepackage{lmodern}
\usepackage[spanish]{babel}

% ── Geometría y espaciado ─────────────────────────────────────────────────────
\usepackage[top=2.5cm, bottom=2.5cm, left=2.8cm, right=2.8cm]{geometry}
\usepackage{setspace}
\setstretch{1.15}
\usepackage{parskip}

% ── Tablas ────────────────────────────────────────────────────────────────────
\usepackage{booktabs}
\usepackage{tabularx}
\usepackage{array}
\usepackage{longtable}
\usepackage{enumitem}

% ── Colores y cajas ───────────────────────────────────────────────────────────
\usepackage[dvipsnames]{xcolor}
\usepackage{tcolorbox}
\tcbuselibrary{skins, breakable}

% ── Cabeceras y pies de página ────────────────────────────────────────────────
\usepackage{fancyhdr}
\usepackage{lastpage}
\pagestyle{fancy}
\fancyhf{}
\fancyhead[L]{\small\textbf{Informe de Evaluación} --- Física para Bioanalistas}
\fancyhead[R]{\small mariana salazar 33570679}
\fancyfoot[C]{\small Página \thepage\ de \pageref{LastPage}}
\renewcommand{\headrulewidth}{0.4pt}

% ── Hiperenlaces ──────────────────────────────────────────────────────────────
\usepackage[colorlinks=true, linkcolor=NavyBlue, urlcolor=NavyBlue]{hyperref}

% ── Paleta de colores personalizados ─────────────────────────────────────────
\definecolor{desempeno_excelente}{HTML}{1A6B2F}
\definecolor{desempeno_bueno}{HTML}{2E5A9C}
\definecolor{desempeno_suficiente}{HTML}{8B6914}
\definecolor{desempeno_deficiente}{HTML}{C0392B}
\definecolor{desempeno_insuficiente}{HTML}{7B241C}
\definecolor{grisFondo}{HTML}{F5F5F5}
\definecolor{azulTitulo}{HTML}{1B2A6B}
\definecolor{verdeFortaleza}{HTML}{1A6B2F}
\definecolor{naranjaMejora}{HTML}{A04000}

% ── Estilos de cajas ─────────────────────────────────────────────────────────
\tcbset{
  cajaTitulo/.style={
    enhanced, breakable,
    colback=azulTitulo!8, colframe=azulTitulo,
    fonttitle=\bfseries\large, coltitle=white,
    attach boxed title to top left={yshift=-2mm, xshift=4mm},
    boxed title style={colback=azulTitulo},
    top=6pt, bottom=6pt, left=8pt, right=8pt,
  },
  cajaFortaleza/.style={
    enhanced, breakable,
    colback=verdeFortaleza!6, colframe=verdeFortaleza!60,
    fonttitle=\bfseries, coltitle=verdeFortaleza,
    top=4pt, bottom=4pt, left=8pt, right=8pt,
  },
  cajaMejora/.style={
    enhanced, breakable,
    colback=naranjaMejora!6, colframe=naranjaMejora!60,
    fonttitle=\bfseries, coltitle=naranjaMejora,
    top=4pt, bottom=4pt, left=8pt, right=8pt,
  },
  cajaRecomendacion/.style={
    enhanced, breakable,
    colback=grisFondo, colframe=gray!50,
    fonttitle=\bfseries, coltitle=black,
    top=4pt, bottom=4pt, left=8pt, right=8pt,
  },
}

% ─────────────────────────────────────────────────────────────────────────────
\begin{document}

% ══ PORTADA ══════════════════════════════════════════════════════════════════
\begin{center}
  {\Large\bfseries\color{azulTitulo} Universidad de Oriente/Núcleo Sucre --- Física para Bioanalistas}\\[4pt]
  {\large Informe de Evaluación Preliminar}\\[12pt]
  \rule{\linewidth}{1.2pt}\\[6pt]
  {\LARGE\bfseries mariana salazar 33570679}\\[4pt]
  \rule{\linewidth}{0.6pt}
\end{center}

% ══ FICHA TÉCNICA ════════════════════════════════════════════════════════════
\vspace{4pt}
\begin{tcolorbox}[cajaTitulo, title=Datos del informe]
\begin{tabular}{@{}llll@{}}
  \textbf{Estudiante:}  & mariana salazar 33570679  &
  \textbf{Fecha:}       & 2026-06-26 \\[4pt]
  \textbf{Archivo PDF:} & \texttt{mariana\_salazar\_33570679.pdf}  &
  \textbf{Páginas:}     & 6 \\
\end{tabular}
\end{tcolorbox}

% ══ RESULTADO GLOBAL ═════════════════════════════════════════════════════════
\vspace{6pt}
\begin{center}
  \begin{tcolorbox}[
    enhanced,
    colback=white, colframe=desempeno_deficiente,
    width=0.62\linewidth,
    boxrule=2pt,
    halign=center, valign=center,
    top=8pt, bottom=8pt,
  ]
    \centering
    {\large\bfseries Puntaje sugerido}\\[6pt]
    {\Huge\bfseries\color{desempeno_deficiente} 5.0/10}\\[6pt]
    {\large Nivel de desempeño: \textbf{\color{desempeno_deficiente} Deficiente}}
  \end{tcolorbox}
\end{center}

% ══ RESUMEN GENERAL ══════════════════════════════════════════════════════════
\section*{Resumen general}
El estudiante demuestra un conocimiento básico de las fórmulas y principios físicos, pero comete errores conceptuales y procedimentales significativos en varias preguntas. Hay problemas recurrentes con la aplicación correcta de las ecuaciones (signos, inclusión de todas las variables relevantes) y errores de cálculo o de unidades que afectan la validez de los resultados finales. La claridad y organización son adecuadas, pero la precisión es una debilidad.

% ══ FORTALEZAS Y ÁREAS DE MEJORA ════════════════════════════════════════════
\vspace{4pt}
\begin{minipage}[t]{0.48\linewidth}
  \begin{tcolorbox}[cajaFortaleza, title={Fortalezas}]
\begin{itemize}[leftmargin=1.5em, itemsep=2pt]
  \item Correcta identificación de las fórmulas principales en la mayoría de los casos.
  \item Buena organización en la presentación de datos y fórmulas.
  \item Correcta aplicación de la Ley de Poiseuille en la Pregunta 5, incluyendo conversiones de unidades y cálculos precisos.
  \item Extracción de datos iniciales generalmente correcta.
\end{itemize}
  \end{tcolorbox}
\end{minipage}
\hfill
\begin{minipage}[t]{0.48\linewidth}
  \begin{tcolorbox}[cajaMejora, title={Areas de mejora}]
\begin{itemize}[leftmargin=1.5em, itemsep=2pt]
  \item Profundizar en la comprensión conceptual de los principios físicos, especialmente en la aplicación de la ecuación de Bernoulli y la Ley de Poiseuille en escenarios complejos.
  \item Mejorar la precisión en los cálculos numéricos y la consistencia en el uso de los datos a lo largo de un problema.
  \item Prestar más atención a las unidades y conversiones, así como a la magnitud de los resultados finales.
  \item Revisar cuidadosamente la lectura de los problemas para identificar todas las condiciones y datos relevantes, especialmente cuando hay múltiples factores en juego.
\end{itemize}
  \end{tcolorbox}
\end{minipage}

% ══ OBSERVACIONES POR PREGUNTA ═══════════════════════════════════════════════
\section*{Observaciones por pregunta}

\begin{longtable}{>{\bfseries}p{2.8cm} p{1.6cm} p{7.0cm} p{2.8cm}}
  \toprule
  \textbf{Pregunta} & \textbf{Puntaje} & \textbf{Evaluación} & \textbf{Errores detectados} \\
  \midrule
  \endhead
  \bottomrule
  \endfoot
    \textbf{Pregunta 1: Presión hidrostática y venosa} & 1.0/2.0 & El estudiante aplica correctamente la fórmula de presión hidrostática (P = $\rho$gh) y realiza el cálculo numérico para obtener la altura en metros (0.24 m). Sin embargo, comete un error significativo al convertir el resultado a centímetros, escribiendo '0,24 cm' en lugar de '24 cm'. Este es un error de magnitud importante. & \textit{Error de conversión de unidades (0.24 m a 0.24 cm).} \\
    \midrule
    \textbf{Pregunta 2: Principio de Arquímedes} & 1.0/2.0 & El estudiante identifica correctamente el principio de Arquímedes y las fórmulas para la fuerza de empuje (E = W\_real - W\_aparente) y la densidad ($\rho$ = m/V). Calcula correctamente la masa real de la muestra. Sin embargo, al calcular el volumen de la muestra a partir de la fuerza de empuje, utiliza un valor incorrecto para la fuerza de empuje (240 N en lugar de 2.40 N), lo que lleva a un volumen incorrecto y, por ende, a una densidad final incorrecta. El procedimiento general es correcto, pero el error numérico es significativo. & \textit{Error procedimental en el cálculo del volumen (uso de valor incorrecto para la fuerza de empuje).} \\
    \midrule
    \textbf{Pregunta 3: Hemodinámica (Continuidad y Bernoulli)} & 0.5/2.0 & El estudiante aplica correctamente la ecuación de continuidad para calcular la velocidad en la estenosis (V$_{2}$ = 1.20 m/s). Sin embargo, al aplicar la ecuación de Bernoulli, comete un error conceptual en el signo del término de energía cinética al despejar la presión final (utiliza P$_{2}$ = P$_{1}$ + ½$\rho$(V$_{2}$$^{2}$ - V$_{1}$$^{2}$) en lugar de P$_{2}$ = P$_{1}$ + ½$\rho$(V$_{1}$$^{2}$ - V$_{2}$$^{2}$)). Además, hay una inconsistencia numérica en el cálculo final, ya que el resultado de 13284.5 Pa no se deriva directamente de los pasos intermedios mostrados (que darían 13325 Pa con su fórmula incorrecta). & \textit{Error conceptual en la aplicación de la ecuación de Bernoulli (signo del término cinético).; Error procedimental en el cálculo final (inconsistencia numérica).} \\
    \midrule
    \textbf{Pregunta 4: Ley de Poiseuille (Análisis Proporcional)} & 0.5/2.0 & El estudiante identifica correctamente la relación de Poiseuille con el radio (Q $\propto$ r$^{4}$) y calcula correctamente el efecto de la reducción del radio (0.8)$^{4}$. Sin embargo, ignora el dato de que la viscosidad de la sangre aumenta en un 50\%, lo cual es un factor crucial en la Ley de Poiseuille. Al no incluir el cambio de viscosidad, la respuesta es conceptualmente incompleta. La pregunta presenta una ambigüedad en su enunciado que podría haber confundido al estudiante. & \textit{Error conceptual al no considerar el cambio de viscosidad en la Ley de Poiseuille.} \\
    \midrule
    \textbf{Pregunta 5: Flujo Viscoso y Resistencia} & 2.0/2.0 & El estudiante extrae correctamente los datos, realiza las conversiones de unidades necesarias (diámetro a radio en metros) y aplica correctamente la Ley de Poiseuille despejada para la diferencia de presión ($\Delta$P). Los cálculos son precisos y el resultado final es correcto, incluyendo las unidades (Pa). & \textit{---} \\
    \midrule
\end{longtable}

% ══ ERRORES DETECTADOS ═══════════════════════════════════════════════════════
\section*{Análisis de errores}

\subsection*{Errores conceptuales}
\begin{itemize}[leftmargin=1.5em, itemsep=2pt]
  \item Malinterpretación de la ecuación de Bernoulli al despejar la presión, resultando en un signo incorrecto para el término de energía cinética (Pregunta 3).
  \item Omisión de una variable clave (viscosidad) en la aplicación de la Ley de Poiseuille, a pesar de que se especificó un cambio en ella (Pregunta 4).
\end{itemize}

\subsection*{Errores procedimentales}
\begin{itemize}[leftmargin=1.5em, itemsep=2pt]
  \item Error de conversión de unidades (metros a centímetros) que altera la magnitud del resultado (Pregunta 1).
  \item Error numérico al transcribir un valor (2.40 N como 240 N) que se propaga en cálculos posteriores (Pregunta 2).
  \item Inconsistencia en el cálculo final de la Pregunta 3, donde el resultado numérico no coincide con los pasos intermedios.
\end{itemize}

% ══ RECOMENDACIONES AL ESTUDIANTE ════════════════════════════════════════════
\section*{Retroalimentación para el estudiante}
\begin{tcolorbox}[cajaRecomendacion, title={Dirigido a: mariana salazar 33570679}]
Mariana, tu examen muestra que tienes una base para entender los principios de la física, pero es crucial que prestes más atención a los detalles. Revisa cuidadosamente la formulación de las ecuaciones, especialmente los signos y la inclusión de todas las variables relevantes, como en la ecuación de Bernoulli y la Ley de Poiseuille. Practica la conversión de unidades y la consistencia en los cálculos para evitar errores de magnitud. Te sugiero repasar los conceptos de fuerza de empuje y cómo se relaciona con el volumen, así como la aplicación completa de la Ley de Poiseuille. Una revisión sistemática de tus pasos y una doble verificación de los resultados te ayudarán a mejorar significativamente.
\end{tcolorbox}

% ══ NOTA PARA EL PROFESOR ════════════════════════════════════════════════════
\section*{Nota para el profesor}
\begin{tcolorbox}[
  enhanced, breakable,
  colback=yellow!8, colframe=yellow!60!black,
  fonttitle=\bfseries, coltitle=black,
  title={Observaciones para revisión manual},
  top=4pt, bottom=4pt, left=8pt, right=8pt,
]
La Pregunta 4 presenta una ambigüedad en su enunciado: primero indica que 'la viscosidad de la sangre se mantienen rigurosamente constantes' y luego, en la misma frase, añade 'y la viscosidad de la sangre aumenta en un 50\%'. El estudiante optó por considerar solo el cambio de radio, ignorando el aumento de viscosidad. Sugiero considerar esta ambigüedad al evaluar la respuesta del estudiante para esta pregunta. He evaluado asumiendo que el aumento del 50\% era la condición a aplicar, lo que implica un error conceptual por parte del estudiante.
\end{tcolorbox}

% ══ PIE DE INFORME ════════════════════════════════════════════════════════════
\vspace{12pt}
\begin{center}
  \footnotesize\color{gray}
  Informe generado automáticamente por el sistema de evaluación con IA (gemini-2.5-flash) y soporte LaTex. \\
  Este documento es una evaluación \textbf{preliminar} para ser revisado por el profesor antes de ser entregado al 
  estudiante. \\
  Generado el 2026-06-26 \textbullet\ BIOANALISIS Sistema Académico de Evaluación.
  \end{center}

\end{document}
